% Independently derived finite-arithmetic appendix, prepared by repo_foundations.
% Intended as an input section, not a standalone document.
\section{Finite contour cancellation and uniform arithmetic schedules}
\label{sec:finite-contour-audit}

This section records concrete replacements for the analytic notation used
in the contour argument. In particular, it does not infer provability from
the mere existence of a primitive recursive approximation algorithm.

\subsection{A strictly finite Cauchy lemma}

Use the taxicab norm $\lVert u+iv\rVert_1=|u|+|v|$ on rational complex
numbers. Let a rational rectangle have width $W$ and height $H$, and fix
an integer $R\ge W+H$. Suppose that, on the rectangle, a function $G$
has rational approximation procedures for $G$ and $G'$ with error at most
$\eta$, and a proved uniform Taylor inequality
\[
 \lVert G(z)-G(a)-G'(a)(z-a)\rVert_1
       \le M\lVert z-a\rVert_1^2
\]
whenever the segment from $a$ to $z$ stays in the rectangle. The displayed
inequality means its finite rational enclosure version with the indicated
evaluation errors. Here $M\ge1$ is an integer bound, obtained from the
concrete formula for $G$; it is not selected by compactness.

Divide each side into $N$ equal parts. Each of the $N^2$ small rectangles
has taxicab diameter $d=(W+H)/N$ and perimeter $2d$. On each cell choose
its lower left vertex $a$ and the rational affine function
\[
 A_a(z)=G_\eta(a)+G'_\eta(a)(z-a).
\]
For every cell vertex $v$ the rational residual satisfies
\[
 \lVert G_\eta(v)-A_a(v)\rVert_1
       \le Md^2+\eta(2+d).
\]
The trapezoidal boundary sum of an affine complex function is exactly
zero: the constant part telescopes, while the linear part uses
$(u+v)(v-u)=v^2-u^2$. Hence the cell boundary sum is at most
$2d[Md^2+\eta(2+d)]$. Every internal edge is used twice with opposite
orientations and the same rational vertex values. It cancels exactly.
The remaining outer boundary trapezoidal sum $C_{N,\eta}$ therefore obeys
\begin{equation}\label{eq:finite-rectangle}
 \lVert C_{N,\eta}\rVert_1
 \le \frac{2MR^3}{N}+4RN\eta+2R^2\eta.
\end{equation}
This proves the finite lemma using only distributivity, the triangle
inequality, finite summation, and exact incidence cancellation. In
particular the evaluation allowance includes the growing internal
perimeter; omitting its factor $N$ would be incorrect.

For requested error $2^{-q}$ take
\[
 N=2^{q+6}(M+1)(R+1)^3,
 \qquad
 \eta\le\frac{2^{-q-8}}{(R+1)^2(N+1)}.
\]
The three terms of \eqref{eq:finite-rectangle} are respectively at most
$2^{-q-5},2^{-q-6},2^{-q-7}$, so their sum is smaller than $2^{-q-4}$.
If $b(u)$ is the least integer with $u\le2^{b(u)}$, a sufficient sample
precision is
\[
 q_{\rm ev}=q+8+2b(R+1)+b(N+1).
\]
The search for $b(u)$ is bounded by $u$. Enlarge $M$, if needed, to also bound $|G'|$. This gives the usual
finite rectangle-refinement estimate for the contour integral name; its
mesh error is at most a constant times $MR^2/N$. Increasing $q$ by four
at the start absorbs this additional comparison. Thus the integral
notation abbreviates a theorem about finite rational tables at every
requested precision, rather than an appeal to a general integration
theorem.

\subsection{A concrete derivative bound after pole subtraction}

For the PNT contour put
\[
 A(s)=\frac{y^{s+1}}{s(s+1)},\qquad
 G(s)=-\frac{\zeta'(s)}{\zeta(s)}A(s)-\frac{y^2}{2(s-1)}.
\]
The removable value at $s=1$ is evaluated by its regular expression below.
One deliberately generous usable bound, in the taxicab norm, is
\[
 M=2^{512}(\lceil y\rceil+1)^8(T+4)^{40}.
\]
It bounds the magnitudes of $G,G',G''$ and therefore its Taylor remainder,
with room for fixed norm-conversion factors. Here is a direct check.

Outside $|s-1|<1/128$, Euler summation with lower cutoff $1$ gives,
for $0\le j\le3$, bounds for $\zeta^{(j)}$ using
\[
 \left|\int_1^\infty\{u\}u^{-s-1}(\log u)^j\,du\right|
       \le\frac{j!}{\sigma^{j+1}}\le j!2^{j+1}.
\]
The pole derivatives are at most $j!128^{j+1}$. Thus each
$|\zeta^{(j)}|$ is bounded by $Z=2^{40}(T+3)$.
The proved strip apartness gives
$|1/\zeta|\le J=2^{71}(T+3)^7$. On the remaining low-ordinate band the same reciprocal bound follows from
$|1/\zeta|=|s-1|/|w|\le(1/32)/(7/8)<1\le J$.
With $U=ZJ$, ordinary quotient
identities bound $|F|,|F'|,|F''|$, for $F=-\zeta'/\zeta$, respectively
by $U,2U^2,6U^3$. The identities
\[
 \frac1{s(s+1)}=\frac1s-\frac1{s+1},\qquad
 \frac{d^j}{ds^j}y^{s+1}=(\log y)^j y^{s+1}
\]
show that $|A^{(j)}|\le2^6 y^6$ for $0\le j\le3$ on
$1/2\le\sigma\le2$, since $\log y\le y$ and $y\ge2$.
Expanding $G''$ by the product rule gives a bound smaller than
$2^{360}(T+3)^{24}y^6$.

Near the pole write $w=(s-1)\zeta=1+(s-1)h$ and
$g=-w'/w$. On the rational band
$|\sigma-1|\le1/64$, $|t|\le1/64$, the previously proved estimates
give $|w|\ge7/8$. Euler's formula bounds the first three derivatives
of $h$ by $5,10,40,240$, respectively. Hence the first three
derivatives of $w$ are bounded by $6,22,128$, and
$|g|,|g'|,|g''|\le2^{13}$. On this band use the stable expression
\[
 G(s)=\int_0^1 A'(1+u(s-1))\,du+g(s)A(s).
\]
Its second derivative is bounded directly by the displayed bounds for
$A^{(j)}$ and $g^{(j)}$; it is much smaller than the proposed $M$.
The contour can be partitioned at the rational lines $t=\pm1/64$.
Thus choosing between the two formulas requires only rational comparisons,
not a decision about an unknown complex zero. The tiny horizontal
distance $\delta_T$ from the pole causes no numerical instability in
the regular expression.

The simple-pole integral is also finite. Excise a rational square centred
at $1$ of half-side $\rho=\tfrac12\min\{\delta_T,c-1,1/64\}$,
which is positive and fits inside both vertical sides; cancel the rational grid between it
and the outer rectangle, and compute its four sides. Their real parts
cancel in pairs and their imaginary parts add to
$8i\int_0^1du/(1+u^2)=2\pi i$, where this identity defines the same
explicit rational integral name of $\pi$ as usual. In particular
$2\le\pi\le4$, directly from the integrand bounds. This proves the
residue calculation without an additional infinitary residue theorem.

\subsection{Fixed rational logarithm and exponential algorithms}

For a rational $u\in[1,2]$, put $z=(u-1)/(u+1)\in[0,1/3]$.
The finite sum
\[
 \ell_N(u)=2\sum_{j=0}^{N-1}\frac{z^{2j+1}}{2j+1}
\]
has tail at most $9^{-N}$, since
$2\sum_{j\ge N}|z|^{2j+1}\le(3/4)9^{-N}$.
For any positive rational $a/b\ge1$, bounded integer search finds
$j$ with $2^j\le a/b<2^{j+1}$. Write $a/b=2^ju$ and approximate
its logarithm by $j\ell_N(2)+\ell_N(u)$. The error is at most
$(j+1)9^{-N}$. For example $N=p+j+2$ gives error smaller than
$2^{-p}$; for $a/b<1$ use its reciprocal and change sign.
The identities for these names follow from finite geometric-series
identities or rational rectangle refinements, with these same tails.

For a rational complex $z$ with $\lVert z\rVert_1\le B$, where $B\ge1$ is an integer,
use the Taylor polynomial $\sum_{j<N}z^j/j!$. At index $2B$ its
absolute coefficient is at most $B^B\le2^{B^2}$, since the last $B$
factors of $(2B)!$ are at least $B$. Every subsequent ratio is at most
$1/2$. Thus
\[
 \sum_{j\ge N}\frac{B^j}{j!}
       \le2^{1+B^2+2B-N}\qquad(N\ge2B).
\]
Taking $N=B^2+2B+p+2$ gives error at most $2^{-p-1}$.
The case $B=0$ is exact. The derivative and multiplication identities
are obtained by differentiating the finite polynomial and applying the
same explicit tails. When the input itself is approximate, the series
gives the Lipschitz bound $4^{B+1}$ on the ball of radius $B$; the
input precision is chosen first to absorb this displayed factor.

For Euler summation and its first three derivatives, truncation at
$V=2^{2v}$ has a bound of the form
\[
 2^{12}(T+3)(v+1)^3 2^{-v}.
\]
Repeated integration by parts gives this explicitly from
\[
 \int_V^\infty u^{-\sigma-1}(\log u)^jdu
 =V^{-\sigma}\sum_{h=0}^j
    \frac{j!}{(j-h)!}\frac{(\log V)^{j-h}}{\sigma^{h+1}}.
\]
For instance $v=4(p+T+32)^2$ is sufficient for error $2^{-p}$.
The inequality $(v+1)^3\le2^{v/2}$ used here holds for $v\ge64$
and has an elementary finite-induction proof. Split the remaining
integral at all integers up to $V$. On each interval its fractional-part
factor is affine, and a crude derivative bound is
$K=2^{16}(T+3)(V+1)^4$. A rational mesh of step $1/Q$ with
$Q\ge2^{p+4}K(V+1)$ therefore has total quadrature error at most
$2^{-p-4}$. Evaluate its exponential and logarithm terms with enough
additional precision to absorb its finite coefficient sum. All cutoffs
are displayed bounded arithmetic constructions.

For a finite Euler product at real part $1+d$, $d>0$ rational, its
distance from the Dirichlet series is at most $P^{-d}/d$ when all
primes up to the integer $P$ are used. Choose
$P=2^V$ with $V\ge(p+d^{-1}+10)/d$. Then this tail is below $2^{-p}$.
Each factor has denominator apartness at least $1/2$. The real
logarithm of its squared modulus may be expanded with the elementary
geometric tail, and the identity
$3+4\cos v+\cos2v=2(1+\cos v)^2$ applies to every finite summand.
Choose each factor's input evaluation tolerance smaller than the desired
final error divided by $8P4^P[2+2/d]^{16}$. Reciprocal evaluation loses
at most a factor $8$, and a product of at most $P$ factors of magnitude
at most $2$ has error at most $P4^P$ times the factor error when all
factor errors are at most $1$. The remaining displayed multiplier
absorbs the eight-factor positivity comparison. This gives the finite version of the Euler-product
positivity inequality with explicit slack.

\subsection{Uniform expansion order and the induction discipline}

The factorial expansion need not be proved by induction on a statement
quantifying over real precisions. For the requested integer order $r$,
form the single finite expression
\[
 S_r(t)=t\sum_{j<r}\frac{j!}{(\log t)^{j+1}}.
\]
The derivative identity
\[
 S'_r(t)=\frac1{\log t}-\frac{r!}{(\log t)^{r+1}}
\]
is exact finite telescoping: the terms with adjacent indices cancel.
On $2\le t\le n$, an integer bound for $|S''_r|$ is
$B_r^*=2^{r+4}(r+2)^4r!$. Telescope over a rational mesh of step
at most $2^{-m}$, with $m=p+B_r^*+n+4$. Its total Taylor error is
at most $B_r^*n2^{-m}<2^{-p-3}$.
On $t\ge2$ clip each rational logarithm sample upward to $1/2$;
the true value is at least $1/2$, so clipping cannot increase its error.
Consequently every sampled reciprocal has the same positive lower
denominator bound used in the derivative estimates.
All logarithm powers and factorials are evaluated by fixed PR programs
with the parameter $r$; no universal evaluator for arbitrary PR codes
is being assumed. Logarithm evaluation tolerances can, for example,
be taken below $2^{-p-10}/[(n+1)(B_r^*+1)]$ to absorb the finite
coefficient derivatives. This yields the expansion identity uniformly
in $r$ after denominators are cleared.

In every application, fix the outer numerical parameters, the requested
precision, all cutoffs, and all rational denominator bounds before
induction. A loop invariant is an equality or inequality between
primitive recursive integer expressions, or the Boolean conjunction of
such checks over the already processed prefix. Quantifier-free induction
proves the row invariant and its finite fold. Universal introduction is
used only after that derivation. In particular, there is no induction on
``all precisions are correct'', ``the function is holomorphic'', or
``all sufficiently large arguments satisfy the estimate''. The analytic
notation is justified by the proved finite identities and their explicit
error budgets, rather than by a generic assertion that computability
implies arithmetic provability.
